\chapter{Accessibility Without Peripheralization}\label{ch:accessibility}

Accessibility in cinema has usually been delivered at the edges: captions at the bottom of the frame, a sign-language interpreter in a corner or a separate screening, audio description through headphones, a device in the cupholder. The accommodation is real and the placement says something. The film is for the typical viewer, and access is attached to it. Selective cinema offers a different arrangement, in which an accessible realization occupies the whole screen as a member of the family. This chapter considers what that makes possible and what it obliges.

\section{Access streams}

Several kinds of accommodation can be carried as streams.

A \emph{sign-language realization} can place the signing performer within the composition rather than beside it: framed as part of the scene, or as a character's inner voice, with the film's images composed around the signing rather than shrunk to make room for it.

A \emph{captioned realization} can carry text designed as part of the image, placed where it reads best in each shot, rather than confined to a fixed band.

A \emph{low-intensity realization} can reduce flashing, rapid cutting, sudden loudness in the picture, and visual clutter for viewers with photosensitivity, vestibular sensitivity, or sensory overload, as an intensity transformation in the sense of \Cref{ch:divergence}.

A \emph{focal realization} can enlarge the action that matters, simplify the composition, and hold shots longer for viewers with low vision or slower visual processing.

What these share is that they are visual. That matters, because \Cref{ch:composite-sound} showed that visual variation fits naturally with shared sound. A signing realization, a captioned realization, and a focal realization can all run under the same complementary score, and their viewers remain part of the room's shared audio. Accommodations that are carried by sound, such as audio description, need individual sound and carry its costs, but even these can be designed with the family rather than added afterward.

\section{Parity}

Carrying access on the main screen does not by itself make it equal. A system that can run several streams can still give the accessible one the least: the weakest writing, the least rehearsed performance, the lowest share of light, the fewest distinct frames. The danger is peripheralization moved from the edge of the frame into the schedule.

The obligation this chapter proposes is \emph{parity}. An accessible realization should be judged as a complete film, by the same standards as the others, not by whether it transmits the minimum information a viewer needs to follow the plot. A signing realization should be composed by people who understand sign language as a visual art, not produced by pasting an interpreter into existing footage. A focal realization should be edited for its viewers, not cropped from someone else's cut.

Parity has a concrete meaning in the terms of Part~I. Every stream in a family receives some share of light and flash rate. Parity requires that access streams receive the same share as the others under the same conditions, and that the family's protection profile and switching apply to them equally. An access stream that is dimmer, lower in frame rate, or excluded from switching is peripheral however central its position on the screen.

\section{Competition for capacity}

Parity creates a hard choice that should be faced openly. \Cref{ch:capacity} showed that a present-day projector can carry two or three streams. If one of them is an access stream, it is not available for a second genre, a second depth, or a second ending. Access competes with artistic variation for the same small budget.

There are three responses. The first is to use the budget efficiently. Access realizations often share most of their frames with another realization: a captioned realization differs from its uncaptioned counterpart only where text appears, and adaptive gating keeps the glasses open everywhere else. Such access streams cost little light. The second is to design access into the family rather than beside it. A family whose realizations are the complementary genre pair of \Cref{ch:genre} could caption both, or make one of them the focal realization, so that access is a property of existing realizations rather than a stream of its own. The third is to decide, and to say, that access comes first when a choice must be made. The book takes that position, on the ground that an artistic dimension that excludes some viewers is worth less than one that includes them.

\section{The glasses themselves}

Selective cinema has an access problem of its own that should not be hidden. The glasses must be worn. Some viewers cannot wear them comfortably, some wear prescription lenses that make fitting harder, and the flicker of the shutters, though well above the rates most associated with photosensitive seizures, is not tolerated by everyone. Adaptive gating adds a further consideration: its brightness steps, if they recurred rhythmically at low frequencies, could produce exactly the kind of flashing that photosensitive viewers must avoid. Editors should never schedule gating changes as a regular pulse, and the protection profile should include limits on how often brightness may change.

For viewers who cannot use the glasses at all, the composite of \Cref{ch:unfiltered} is what they see. In coincident passages it is the film, correct and at full brightness. In divergent passages it is whatever the family made of it. A family that designs its composite gives these viewers a real realization; one that ignores it gives them a double exposure. Designing the composite is therefore partly an access decision, and a family that cares about access should treat the unaided view as a realization owed to the people who have no other.

\section{Access as a design dimension}

The chapter's argument reverses the usual order. Accessibility is usually something done to a finished film. In a variant family it can be one of the dimensions the family is built on from the start, with its own column in the variant matrix, its own checks, and its own claim on the budget. That is more demanding than adding captions afterward, and it is the only way to give access the central position selective cinema makes available rather than rebuilding the old periphery inside the new schedule.
